% Build: latexmk -pdf project1_solution.tex
\documentclass[a4paper,10pt]{article}
\usepackage[a4paper,left=8mm,right=8mm,top=10mm,bottom=20mm]{geometry}
\usepackage[utf8]{inputenc}
\usepackage{kotex}
\usepackage{amsmath,amsfonts,amssymb,mathrsfs,mathtools}
\usepackage{amsthm,stmaryrd,enumitem,array,booktabs}
\usepackage[hidelinks]{hyperref}
\title{프로그램분석 26F Project 1}
\author{2024-21504 남연우(Yeonwoo Nam)}
\date{}
\newtheorem{theorem}{정리}
\newtheorem{lemma}{렘마}
\newcommand{\ps}[1]{\mathcal P(#1)}
\newcommand{\lf}{\operatorname{lfp}}
\newcommand{\lab}{\operatorname{label}}
\newcommand{\halt}{\mathsf{halt}}
\newcommand{\nxt}{\operatorname{next}}
\newcommand{\nt}{\operatorname{nextTrue}}
\newcommand{\nf}{\operatorname{nextFalse}}
\newcommand{\cont}{\operatorname{cont}}
\newcommand{\Eval}{\operatorname{Eval}}
\newcommand{\Id}{\operatorname{id}}
\newcommand{\tr}{\hookrightarrow}
\newcommand{\atr}{\hookrightarrow^\sharp}
\newcommand{\T}{\mathsf T}
\newcommand{\Ff}{\mathsf F}
\newcommand{\addr}{\operatorname{addr}}
\newcommand{\Step}{\operatorname{Step}}
\newcommand{\Graph}{\operatorname{Graph}}
\begin{document}
\input{macros.tex}
\maketitle

\paragraph{공통 설정.}
분석할 프로그램 $P$를 고정하고, 그 안의 모든 labeled command의 집합을
$\mathrm{Cmd}_P$, 원래 레이블의 유한집합을 $L_P\subseteq\mathbb N$이라 하자.
각 command에는 서로 다른 레이블이 있으므로 $\lab:\mathrm{Cmd}_P\to L_P$는 전단사이다.
$P_l$은 레이블 $l$의 command이다.
종료를 나타내는 새 기호 $\halt\notin\mathbb N$을 추가하여
$L=L_P\cup\{\halt\}$로 둔다. $\halt$는 command나 \texttt{goto}의 대상이 아니다.

프로그램의 변수집합 $V_P$와 단사함수 $\addr:V_P\to\mathbb Z$를 고정한다.
값, 메모리, 모둠 메모리 공간은 각각
\[
 \mathrm{Val}=\mathbb Z,\qquad M=V_P\to\mathbb Z,\qquad D=\ps M
\]
이다. 주소도 정수값으로 표현하며, $\&x$의 값은 $\addr(x)$이다.
초기 메모리집합 $M_0\subseteq M$는 주어진 매개변수로 두고,
초기값이 모두 $0$이라는 가정은 하지 않는다.
잘못된 주소를 역참조하거나 존재하지 않는 레이블로 점프하면
그 선택에 따른 실행은 stuck되어 후속 상태 및 정상 종료 결과를 만들지 않는다.
다른 유효한 비결정적 선택이 있다면 그 실행은 계속된다.

조건식에서 $0$은 거짓, $0$이 아닌 정수는 참이다.
\textbf{$\mathsf{repeat}\ C\ E$는 $\mathsf{while}\ E\ C$와 같은 의미로 정의한다.}
즉, 본문 실행 전에 조건을 검사한다.
$\mathsf{readInt}$는 평가할 때마다 독립적으로 임의의 정수를 반환한다.
식은 메모리를 변경하지 않는다.
구조 내부로 점프했을 때에는 그 위치의 정적인 문맥에 따라 계속 실행하며,
점프한 위치에서 원래의 \texttt{goto} 다음으로 되돌아오지 않는다.

\paragraph{식의 의미.}
각 식 $E$에 대해 $\Eval_E:M\to\ps{\mathbb Z}$를 다음과 같이 정의한다.
이는 하나의 메모리에서 가능한 값들을 모두 모은다.
\[
\begin{aligned}
 \Eval_n(m)&=\{n\},&
 \Eval_x(m)&=\{m(x)\},&
 \Eval_{\&x}(m)&=\{\addr(x)\},\\
 \Eval_{*x}(m)&=\{m(y)\mid y\in V_P,\ m(x)=\addr(y)\},&
 \Eval_{\mathsf{readInt}}(m)&=\mathbb Z,&
 \Eval_{-E}(m)&=\{-v\mid v\in\Eval_E(m)\},\\
 \Eval_{E_1+E_2}(m)&=\{v_1+v_2\mid v_i\in\Eval_{E_i}(m)\},\\
 \Eval_{E_1*E_2}(m)&=\{v_1v_2\mid v_i\in\Eval_{E_i}(m)\},\\
 \Eval_{E_1<E_2}(m)&=\{\mathbf1[v_1<v_2]\mid v_i\in\Eval_{E_i}(m)\}.
\end{aligned}
\]
여기서 $\mathbf1[Q]$는 명제 $Q$가 참이면 $1$, 아니면 $0$이다.
잘못된 역참조의 결과는 공집합이다.
같은 식에 등장하는 서로 다른 $\mathsf{readInt}$는 서로 다른 값을 선택할 수 있다.
집합 입력에 대한 자연스러운 확장은
$\mathrm{eval}_E(X)=\bigcup_{m\in X}\Eval_E(m):D\to\ps{\mathbb Z}$이다.

\paragraph{모둠 의미 연산자.}
$b\in\{\T,\Ff\}$에 대해 $B_\T=\mathbb Z\setminus\{0\}$, $B_\Ff=\{0\}$으로 두고,
다음 연산자들을 정의한다.
\[
\begin{aligned}
 A_{x,E}(X)&=\{m[x\mapsto v]\mid m\in X,\ v\in\Eval_E(m)\},\\
 W_{x,E}(X)&=\{m[y\mapsto v]\mid m\in X,\ y\in V_P,\
                 m(x)=\addr(y),\ v\in\Eval_E(m)\},\\
 B_{E,b}(X)&=\{m\in X\mid \Eval_E(m)\cap B_b\ne\emptyset\},\\
 G_{E,k}(X)&=\{m\in X\mid k\in\Eval_E(m)\}\qquad(k\in L_P).
\end{aligned}                                                    \tag{1}
\]
$A,W,B,G:D\to D$이며 $\Id_D:D\to D$와
$\cup:D\times D\to D$도 사용한다.
각 연산자는 임의의 합집합을 보존한다.
특히 대입에서는 좌변의 주소와 우변의 값을 모두 \emph{동일한 대입 전 메모리}에서 계산한다.
$\{m[x\mapsto v]\mid m\in X,v\in\mathrm{eval}_E(X)\}$로 정의하면
서로 다른 메모리의 값이 섞이므로 정확한 모둠의미가 아니다.

\paragraph{요약 의미공간과 연산자 가정.}
구체적인 요약공간의 구현은 정하지 않는다.
완전격자 $M^\sharp$와 다음 Galois connection만 가정한다.
\[
 \alpha_M:D\rightleftarrows M^\sharp:\gamma_M,\qquad
 \alpha_M(X)\sqsubseteq a\ \Longleftrightarrow\
 X\subseteq\gamma_M(a).
                                                               \tag{2}
\]
사용하는 요약 연산자들의 타입은 다음과 같다.
\[
\begin{array}{c|c|c}
 \text{연산자}&\text{구체 타입}&\text{요약 타입}\\ \hline
 A_{x,E},W_{x,E},B_{E,b},G_{E,k}
   &D\to D&M^\sharp\to M^\sharp\\
 \Id_D&D\to D&M^\sharp\to M^\sharp\\
 \cup&D\times D\to D&M^\sharp\times M^\sharp\to M^\sharp
\end{array}
\]
요약 버전은 각각 $A^\sharp,W^\sharp,B^\sharp,G^\sharp$,
$\Id_{M^\sharp}$, $\sqcup$이다.
$A^\sharp,W^\sharp,B^\sharp,G^\sharp$는 단조이고, 각 대응 연산자 $T$에 대해
\[
 T\circ\gamma_M\ \subseteq\ \gamma_M\circ T^\sharp                 \tag{3}
\]
를 만족한다고 가정한다.
또한 $\gamma_M(a)\cup\gamma_M(b)\subseteq\gamma_M(a\sqcup b)$이다.
이는 $\gamma_M$의 단조성에서 따른다.
여기서는 대입과 조건 제한을 메모리 집합 위의 연산자로 직접 요약하므로,
요약 메모리를 변수별 요약값의 곱공간으로 제한하지 않는다.
별도의 식 분석기 $\mathrm{eval}_E^\sharp:M^\sharp\to\mathrm{Val}^\sharp$를 사용하려면
$\ps{\mathbb Z}\rightleftarrows\mathrm{Val}^\sharp$와 그 안전성도 추가하면 되지만,
아래 정의에는 필요하지 않다.
Galois connection만으로 $\gamma_M(\bot^\sharp)=\emptyset$이 따르지는 않으며,
그 성질은 사용하지 않는다.

\section*{Exercise 1. Transitional style}

\paragraph{1. 정적인 후속 관계.}
각 command $C$의 정상 종료 후 이동할 위치를 $\cont(\lab(C))$라 하자.
다음 순회 $\mathrm{build}(C,k)$는 프로그램 문법만 보고 이를 계산한다.
먼저 $\cont(\lab(C))=k$로 두고, $C$의 형태에 따라 다음을 수행한다.
\[
\begin{array}{c|l}
 C&\text{재귀 호출}\\ \hline
 l:(x:=E),\ l:(*x:=E),\ l:(\mathsf{goto}\ E)&\text{없음}\\
 l:(C_1;C_2)&\mathrm{build}(C_1,\lab(C_2)),\
                         \mathrm{build}(C_2,k)\\
 l:(\mathsf{if}\ E\ C_1\ C_2)&
                         \mathrm{build}(C_1,k),\ \mathrm{build}(C_2,k)\\
 l:(\mathsf{repeat}\ C_0\ E)&\mathrm{build}(C_0,l)
\end{array}
\]
$\mathrm{build}(P,\halt)$를 실행한 후 후속 관계를 다음과 같이 정의한다.
표의 --는 해당 관계가 정의되지 않음을 뜻한다.
\[
\begin{array}{c|c|c|c}
 P_l&\nxt(l)&\nt(l)&\nf(l)\\ \hline
 x:=E&\cont(l)&--&--\\
 *x:=E&\cont(l)&--&--\\
 C_1;C_2&\lab(C_1)&--&--\\
 \mathsf{if}\ E\ C_1\ C_2&--&\lab(C_1)&\lab(C_2)\\
 \mathsf{repeat}\ C_0\ E&--&\lab(C_0)&\cont(l)\\
 \mathsf{goto}\ E&--&--&--\\
 \halt&--&--&--
\end{array}                                                       \tag{4}
\]
각 관계는 필요한 위치에서만 단일값을 갖는다.
$\cont(l)$은 문맥 정보이며, 특히 \texttt{goto}의 fall-through 간선이 아니다.
\texttt{goto}의 실제 도착지는 메모리에 따라 정해지므로 $\nxt$로 미리 고정하지 않는다.

예를 들어 다음처럼 모든 command에 레이블을 붙이자.
\[
\begin{split}
 10:\bigl(&20:(\mathsf{repeat}\
     30:(\mathsf{if}\ y\
          40:(x:=x+(-1))\
          50:(\mathsf{goto}\ 70))\ x);\\[-2mm]
     &60:(70:(y:=1);\ 80:(z:=2))\bigr).
\end{split}
\]
정적인 간선은
$10\to20$, $20\xrightarrow{\T}30$, $20\xrightarrow{\Ff}60$,
$30\xrightarrow{\T}40$, $30\xrightarrow{\Ff}50$,
$40\to20$, $60\to70$, $70\to80$, $80\to\halt$이다.
$50$에서 실행되는 점프의 대상은 $70$이며, 이후 $80,\halt$로 진행한다.
$50$으로 되돌아오는 간선은 없다.
예를 들어 $x=1,y=0$인 초기 메모리에서는
$10,20,30,50,70,80,\halt$의 순서로 실행된다.

\paragraph{2. 구체적인 한걸음 관계와 모둠의미.}
$S=L\times M$에 대해 $\tr\subseteq S\times S$를 다음과 같이 정의한다.
\[
\begin{array}{c|l}
 P_l&\text{전이}\\ \hline
 x:=E&(l,m)\tr(\nxt(l),m[x\mapsto v])
                    \quad(v\in\Eval_E(m))\\
 *x:=E&(l,m)\tr(\nxt(l),m[y\mapsto v])
                    \quad(m(x)=\addr(y),\ v\in\Eval_E(m))\\
 C_1;C_2&(l,m)\tr(\nxt(l),m)\\
 \mathsf{if}\ E\ C_1\ C_2&
 (l,m)\tr(\nt(l),m)\quad(\Eval_E(m)\cap B_\T\ne\emptyset)\\
 & (l,m)\tr(\nf(l),m)\quad(\Eval_E(m)\cap B_\Ff\ne\emptyset)\\
 \mathsf{repeat}\ C_0\ E&
 (l,m)\tr(\nt(l),m)\quad(\Eval_E(m)\cap B_\T\ne\emptyset)\\
 & (l,m)\tr(\nf(l),m)\quad(\Eval_E(m)\cap B_\Ff\ne\emptyset)\\
 \mathsf{goto}\ E&(l,m)\tr(k,m)\quad(k\in L_P\cap\Eval_E(m)).
\end{array}                                                       \tag{5}
\]
$\halt$에서는 전이가 없다.
간선별 메모리 전달 연산자를 $T_{lk}:D\to D$라 쓰자.
대입은 각각 $A,W$, sequence는 $\Id_D$, 참/거짓 간선은 $B_{E,\T},B_{E,\Ff}$,
점프의 $k$ 간선은 $G_{E,k}$이다.
같은 도착지에 여러 간선이 있으면 그 기여를 모두 합집합한다.
이 표현으로 (5)는
\[
 \Step(X)_k=\bigcup_{l\to k}T_{lk}(X_l),\qquad
 X_l=\{m\mid(l,m)\in X\}                                         \tag{6}
\]
와 같다. 여기서 점프의 후보 간선은 모든 $k\in L_P$에 대해 만들며,
$G_{E,k}$가 실제 가능한 메모리만 선택한다.

$l_0=\lab(P)$에 대해
\[
 I=\{(l_0,m)\mid m\in M_0\},\qquad
 F:\ps S\to\ps S,\quad F(X)=I\cup\Step(X),\qquad
 \mathcal C=\lf F.                                               \tag{7}
\]
$\Step$은 관계의 직접상(direct image)이므로 임의의 합집합을 보존하고,
$F$는 단조이며 Scott 연속이다. 따라서
\[
 \mathcal C=\bigcup_{n\ge0}F^n(\emptyset)
            =\bigcup_{n\ge0}\Step^n(I).
\]
이는 유한한 실행 접두사에 등장하는 모든 상태를 모으므로,
종료하는 실행뿐 아니라 무한 실행 중 도달하는 상태도 포함한다.

\paragraph{3. 요약 한걸음 관계와 요약의미.}
요약상태 공간은 $\mathcal A=L\to M^\sharp$이며 순서와 join은 점별로 정의한다.
(2)를 레이블별로 올리면 다음 Galois connection을 얻는다.
\[
\begin{aligned}
 \alpha_S(X)(l)&=\alpha_M(X_l),\\
 \gamma_S(a)&=\{(l,m)\mid m\in\gamma_M(a(l))\},\\
 \alpha_S(X)\sqsubseteq a&\ \Longleftrightarrow\ X\subseteq\gamma_S(a).
\end{aligned}                                                     \tag{8}
\]
우선 local pair의 관계
$\atr_{\mathrm{loc}}\subseteq(L\times M^\sharp)\times(L\times M^\sharp)$를
\[
 (l,a)\atr_{\mathrm{loc}}(k,T_{lk}^\sharp(a))\qquad(l\to k)          \tag{9}
\]
로 정의한다. 즉 (5)의 각 규칙에서 $A,W,\Id,B,G$를
$A^\sharp,W^\sharp,\Id,B^\sharp,G^\sharp$로 바꾼다.
점프는 모든 원래 레이블 $k\in L_P$에 대해
$(l,a)\atr_{\mathrm{loc}}(k,G_{E,k}^\sharp(a))$를 만든다.
참/거짓 및 점프 대상이 불가능함을 요약공간이 확인하지 못하면 여러 간선에
동시에 전파해도 안전하다.

이 관계를 모아서 실제 label map의 변환을 정의한다.
\[
\begin{aligned}
 \Step^\sharp&:\mathcal A\to\mathcal A,&
 \Step^\sharp(a)(k)&=\bigsqcup_{l\to k}T_{lk}^\sharp(a(l)),\\
 I^\sharp&=\alpha_S(I),&
 F^\sharp(a)&=I^\sharp\sqcup\Step^\sharp(a),&
 \mathcal C^\sharp&=\lf F^\sharp .
\end{aligned}                                                     \tag{10}
\]
빈 join은 $\bot^\sharp$이다.
$M^\sharp$가 완전격자이고 모든 전달함수가 단조이므로
$\mathcal A$도 완전격자이며 $F^\sharp$의 최소고정점이 존재한다.
일반적인 요약공간에서 유한 반복으로 계산된다는 주장은 하지 않는다.

과제의 관계 타입과 partition/join 표기를 구별하면 다음과 같다.
$Q=L\times M^\sharp$로 두고,
\[
\begin{aligned}
 \Graph &: \mathcal A\to\ps Q,&
 \Graph(a)&=\{(l,a(l))\mid l\in L\},\\
 \mathrm{Post}^\sharp&:\ps Q\to\ps Q,&
 \mathrm{Post}^\sharp(U)&=\{q'\mid \exists q\in U.\ q\atr_{\rm loc}q'\},\\
 \pi&:\ps Q\to(L\to\ps{M^\sharp}),&
 \pi(U)(k)&=\{b\mid(k,b)\in U\},\\
 J&:(L\to\ps{M^\sharp})\to\mathcal A,&
 J(B)(k)&=\bigsqcup B(k).
\end{aligned}
\]
그러면 $\Step^\sharp=J\circ\pi\circ\mathrm{Post}^\sharp\circ\Graph$이다.
이는 강의의 $\mathcal P(\Id,\sqcup)\circ\pi\circ
\breve{\mathcal P}(\atr_{\rm loc})$를 타입까지 명시한 것이다.
과제에 적힌 map-to-map 관계 $\atr\subseteq\mathcal A\times\mathcal A$는
$a\atr a'$ iff $a'=\Step^\sharp(a)$로 정의할 수 있다.
이 map 관계와 partition 이전의 local 관계를 같은 타입으로 혼동하지 않는다.

\paragraph{4--5. 안전성.}
\begin{lemma}[Local edge coverage]
$m\in\gamma_M(a)$이고 $(l,m)\tr(k,m')$이면,
해당 간선 $l\to k$에 대해 $m'\in\gamma_M(T_{lk}^\sharp(a))$이다.
\end{lemma}
\begin{proof}
대입이면 $m'\in A_{x,E}(\gamma_M(a))$이므로 (3)에 의해 성립한다.
간접대입도 주소와 값이 모두 $m$에서 계산된다.\par
따라서 $m'\in W_{x,E}(\gamma_M(a))$이고 (3)을 적용한다.
Sequence이면 $m'=m$이고 전달함수가 항등함수이다.
조건문과 repeat의 참 간선이면
$m\in B_{E,\T}(\gamma_M(a))$이며, 거짓 간선도 $B_{E,\Ff}$에 대해 같다.
마지막으로 점프이면 $k\in L_P\cap\Eval_E(m)$이므로
$m'=m\in G_{E,k}(\gamma_M(a))$이고 $G^\sharp$의 안전성으로 결론을 얻는다.
$\halt$에서는 검사할 전이가 없다.
\end{proof}

관계의 powerset lifting 자체에 대한 과제의 조건도 명확히 쓸 수 있다.
$\Gamma:\ps Q\to\ps S$를
$\Gamma(U)=\{(l,m)\mid\exists(l,a)\in U.\ m\in\gamma_M(a)\}$로 두면,
위 렘마에서
\[
 \Step\circ\Gamma\ \subseteq\
 \Gamma\circ\mathrm{Post}^\sharp                                \tag{11}
\]
를 얻는다. 이는 local 한걸음 관계의
$\breve{\mathcal P}(\tr)\circ\gamma
 \subseteq\gamma\circ\breve{\mathcal P}(\atr_{\rm loc})$
조건이다. Join은 포함관계를 더 넓힐 뿐이다.

\begin{theorem}[Transitional analysis의 안전성]
모든 $a\in\mathcal A$에 대해
$\Step(\gamma_S(a))\subseteq\gamma_S(\Step^\sharp(a))$이며,
\[
             \mathcal C\subseteq\gamma_S(\mathcal C^\sharp).
\]
\end{theorem}
\begin{proof}
$(k,m')\in\Step(\gamma_S(a))$를 잡으면 어떤 $l,m$에 대해
$m\in\gamma_M(a(l))$, $(l,m)\tr(k,m')$이다.
렘마와 join의 상한 성질로
\[
 m'\in\gamma_M(T_{lk}^\sharp(a(l)))
       \subseteq\gamma_M(\Step^\sharp(a)(k)).
\]
또 $I\subseteq\gamma_S(I^\sharp)$이므로
$F(\gamma_S(a))\subseteq\gamma_S(F^\sharp(a))$이다.
$a_*=\lf F^\sharp$라 두면
$F(\gamma_S(a_*))\subseteq\gamma_S(F^\sharp(a_*))=\gamma_S(a_*)$이다.
따라서 $\gamma_S(a_*)$는 $F$의 pre-fixed point이고,
최소고정점의 최소성에 의해 $\lf F\subseteq\gamma_S(a_*)$이다.
이 증명은 $\gamma_M$가 join을 보존한다거나 bottom이 공집합을 뜻한다고 가정하지 않는다.
\end{proof}

\section*{Exercise 2. Definitional interpreter}

\paragraph{1. 실행기의 문맥과 모둠 실행기.}
\texttt{goto}가 있으므로 command의 의미는 고정된 전체 프로그램 $P$의 문맥에 의존한다.
여기서 $\mathcal I_P(C)$는 $C$에 진입한 후 \emph{$P$의 종료점까지}
실행하여 얻는 최종 메모리들의 집합을 반환한다.
따라서 내부 command를 독립적인 부분 프로그램으로 해석하지 않는다.
이 해석에서
$\mathcal I_P(C_1;C_2)=\mathcal I_P(C_2)\circ\mathcal I_P(C_1)$로 놓으면
$C_2$를 중복 실행할 수 있으므로 그 식을 사용하지 않는다.

$\operatorname{Mon}(D,D)$는 포함관계에 대한 단조함수들의 완전격자이며,
순서와 모든 join은 점별이다.
\[
 \mathcal H=\mathrm{Cmd}_P\to\operatorname{Mon}(D,D).
\]
$H\in\mathcal H$에 대해 종료점 호출을 별도로
\[
 \widehat H(k)=
 \begin{cases}H(P_k)&k\in L_P,\\ \Id_D&k=\halt\end{cases}
                                                                    \tag{12}
\]
로 정의한다. 따라서 $\halt$를 새 command로 추가할 필요가 없다.
실행기 방정식의 함수 $\Phi:\mathcal H\to\mathcal H$는 다음과 같다.
\[
\begin{array}{c|l}
 C=P_l&(\Phi H)(C)(X)\\ \hline
 x:=E&\widehat H(\nxt(l))(A_{x,E}(X))\\
 *x:=E&\widehat H(\nxt(l))(W_{x,E}(X))\\
 C_1;C_2&H(C_1)(X)\\
 \mathsf{if}\ E\ C_1\ C_2&
 H(C_1)(B_{E,\T}(X))\ \cup\ H(C_2)(B_{E,\Ff}(X))\\
 \mathsf{repeat}\ C_0\ E&
 H(C_0)(B_{E,\T}(X))\ \cup\
            \widehat H(\cont(l))(B_{E,\Ff}(X))\\
 \mathsf{goto}\ E&
 \displaystyle\bigcup_{k\in L_P}H(P_k)(G_{E,k}(X)).
\end{array}                                                        \tag{13}
\]
본문 $C_0$의 정상 후속 레이블은 (4)에 의해 다시 $l$이다.
모든 의미 연산자의 타입은 (1)의 $D\to D$,
$\Id_D:D\to D$, $\cup:D\times D\to D$이고,
재귀 호출 $H(C)$ 역시 $D\to D$이다.
식의 보조 의미 $\Eval_E$와 $\mathrm{eval}_E$의 타입은 공통 설정에서 정의하였다.

이는 과제에서 요구하는 command-indexed 함수 공간을 그대로 사용한다.
전체 문법 $\mathrm{Cmd}$에 대한 타입
$\mathrm{Cmd}\to D\to D$로 쓰고 싶으면,
$P$에 없는 command에서는 $\Phi(H)(C)=\lambda X.\emptyset$으로 확장하면 된다.
그 좌표는 $P$의 결과에 영향을 주지 않는다.

\begin{lemma}[실행기 최소고정점의 존재와 의미]
$\Phi$는 Scott 연속이고
\[
 \mathcal I_P=\lf\Phi=\bigsqcup_{n\ge0}\Phi^n(\bot_{\mathcal H}),
 \qquad \bot_{\mathcal H}(C)(X)=\emptyset.                         \tag{14}
\]
또한 모든 $C,X$에 대해
\[
 \mathcal I_P(C)(X)=
 \{m'\mid \exists m\in X.\ (\lab(C),m)\tr^*(\halt,m')\}.             \tag{15}
\]
\end{lemma}
\begin{proof}
(13)의 각 입력 전달함수는 고정되어 있고 단조이다.
재귀 환경 $H$는 그 고정된 입력에 적용되는 형태로만 등장한다.
따라서 증가하는 directed family $\{H_i\}$에 대해 각 호출에서
$(\bigsqcup_iH_i)(C)(Y)=\bigcup_iH_i(C)(Y)$이며,
이를 (13)에 대입하면
$\Phi(\bigsqcup_iH_i)=\bigsqcup_i\Phi(H_i)$이다.
종료점의 항등함수는 환경과 무관한 상수이다.
그러므로 Kleene 정리를 적용할 수 있다.

$H_n=\Phi^n(\bot_{\mathcal H})$라 하자.
$n\ge1$일 때 $H_n(C)(X)$는 $\lab(C)$에서 시작하여
최대 $n$번 전이한 뒤 처음 $\halt$에 도착하는 경로의 최종 메모리 집합이다.
이는 $n$에 대한 귀납법으로 보인다.
첫 단계의 전달함수는 (5)의 첫 전이를 정확히 나타내고,
후속 위치가 command이면 $H_{n-1}$이 남은 경로를 처리한다.
후속 위치가 $\halt$이면 (12)의 항등함수가 즉시 결과를 반환한다.
$n=0$에서는 어떤 command의 결과도 없다.
이를 모든 $n$에 대해 합집합하면 (15)이다.
\end{proof}

따라서 주어진 전체 프로그램의 모둠의미는
$\mathcal C_P=(\lf\Phi)(P):D\to D$이다.
종료하지 않는 실행은 최종 메모리를 만들지 않는다.
이는 Exercise 1의 중간 도달 상태를 모두 모으는 의미와 구별된다.

\paragraph{2--3. 요약 실행기.}
요약 실행기 공간과 종료점 호출을
\[
 \mathcal H^\sharp=\mathrm{Cmd}_P\to
       \operatorname{Mon}(M^\sharp,M^\sharp),\qquad
 \widehat K(k)=
 \begin{cases}K(P_k)&k\in L_P,\\ \Id_{M^\sharp}&k=\halt\end{cases}
                                                                    \tag{16}
\]
로 정의한다.
$\Phi^\sharp:\mathcal H^\sharp\to\mathcal H^\sharp$는
(13)과 같은 모양으로 정의한다.
\[
\begin{array}{c|l}
 C=P_l&(\Phi^\sharp K)(C)(a)\\ \hline
 x:=E&\widehat K(\nxt(l))(A^\sharp_{x,E}(a))\\
 *x:=E&\widehat K(\nxt(l))(W^\sharp_{x,E}(a))\\
 C_1;C_2&K(C_1)(a)\\
 \mathsf{if}\ E\ C_1\ C_2&
 K(C_1)(B^\sharp_{E,\T}(a))\ \sqcup\
 K(C_2)(B^\sharp_{E,\Ff}(a))\\
 \mathsf{repeat}\ C_0\ E&
 K(C_0)(B^\sharp_{E,\T}(a))\ \sqcup\
        \widehat K(\cont(l))(B^\sharp_{E,\Ff}(a))\\
 \mathsf{goto}\ E&
 \displaystyle\bigsqcup_{k\in L_P}K(P_k)(G^\sharp_{E,k}(a)).
\end{array}                                                        \tag{17}
\]
대입, store, guard, target 연산자 및 모든 재귀 호출의 요약 입출력 타입은
$M^\sharp\to M^\sharp$이며,
join의 타입은 $M^\sharp\times M^\sharp\to M^\sharp$이다.
종료점의 요약 항등함수도 같은 타입이다.

$\mathcal H^\sharp$는 점별 순서의 완전격자이다.
$\Phi^\sharp$는 환경 인자에 대해 Scott 연속이다.
실제로 (17)의 재귀 호출 입력은 $K$와 무관하게 고정되고,
유한 join과 directed supremum은 서로 교환된다.
따라서 여기서는 개별 $T^\sharp$의 Scott 연속성까지 요구하지 않고
단조성만으로 충분하다.
\[
 \mathcal I_P^\sharp=\lf\Phi^\sharp,\qquad
 \mathcal C_P^\sharp=(\lf\Phi^\sharp)(P):
                  M^\sharp\to M^\sharp.                         \tag{18}
\]
이는 과제의 $F,F^\sharp$를 각각 $\Phi,\Phi^\sharp$로 쓴 것이다.
최소고정점이 존재한다는 사실은 유한 번에 계산할 수 있다는 뜻은 아니다.
이 단계에서는 구체적인 도메인이나 widening을 정하지 않는다.

\paragraph{4--5. 실행기의 안전성.}
환경 사이의 관계 $H\mathrel{\mathcal R}K$를
\[
 \forall C\in\mathrm{Cmd}_P,\ \forall a\in M^\sharp.\quad
 H(C)(\gamma_M(a))\subseteq\gamma_M(K(C)(a))                       \tag{19}
\]
로 정의한다. 종료점의 두 항등함수도 같은 관계를 만족한다.

\begin{lemma}[환경 방정식의 simulation]
$H\mathrel{\mathcal R}K$이면
$\Phi(H)\mathrel{\mathcal R}\Phi^\sharp(K)$이다.
\end{lemma}
\begin{proof}
임의의 간선 전달함수 $T$와 후속 위치 $k$에 대해,
(3), $\widehat H(k)$의 단조성, (19)를 차례로 적용하면
\[
\begin{aligned}
 \widehat H(k)(T(\gamma_M(a)))
 &\subseteq \widehat H(k)(\gamma_M(T^\sharp(a)))\\
 &\subseteq \gamma_M(\widehat K(k)(T^\sharp(a))).
\end{aligned}                                                     \tag{20}
\]
$k=\halt$일 때 두 번째 포함관계는 항등함수의 경우이다.
직접대입과 간접대입은 각각 $T=A,W$로 (20)을 적용한다.
Sequence는 (19)를 첫 command에 직접 적용한다.
조건문은 $T=B_{E,\T},B_{E,\Ff}$에 (20)을 각각 적용하고,
$\gamma_M(u)\cup\gamma_M(v)\subseteq\gamma_M(u\sqcup v)$를 사용한다.
Repeat도 참 가지의 본문 호출과 거짓 가지의 후속 호출에 같은 논증을 적용한다.
Goto는 각 $k\in L_P$에 대해 $T=G_{E,k}$로 (20)을 적용한 후 합집합한다.
모든 요약 기여는 최종 join 이하이므로 $\gamma_M$의 단조성에 의해 결론을 얻는다.
\end{proof}

\begin{theorem}[Definitional interpreter의 안전성]
모든 프로그램 $P$와 $a\in M^\sharp$에 대해
\[
 \mathcal C_P(\gamma_M(a))
       \subseteq\gamma_M(\mathcal C_P^\sharp(a)).
                                                                    \tag{21}
\]
따라서
$\mathcal C_P(M_0)\subseteq
 \gamma_M(\mathcal C_P^\sharp(\alpha_M(M_0)))$이다.
\end{theorem}
\begin{proof}
$K_*=\lf\Phi^\sharp$라 두자.
$\bot_{\mathcal H}\mathrel{\mathcal R}K_*$는 공집합의 포함관계로 성립한다.
Simulation 렘마와 $\Phi^\sharp(K_*)=K_*$를 사용한 귀납법으로
$\Phi^n(\bot_{\mathcal H})\mathrel{\mathcal R}K_*$가 모든 $n$에 대해 성립한다.
그러므로 모든 $C,a$에 대해
\[
 \mathcal I_P(C)(\gamma_M(a))
 =\bigcup_{n\ge0}\Phi^n(\bot_{\mathcal H})(C)(\gamma_M(a))
 \subseteq\gamma_M(K_*(C)(a)).
\]
$C=P$를 대입하면 (21)이다.
마지막 결론은 $M_0\subseteq\gamma_M(\alpha_M(M_0))$ 및
$\mathcal C_P$의 단조성에서 따른다.
이 증명에서는 $\gamma_M$와 무한 join을 교환하지 않는다.
\end{proof}

\paragraph{두 의미의 일치.}
(15)에 의해 같은 초기집합 $M_0$에 대해
\[
 \mathcal C_P(M_0)=
 \{m\mid(\halt,m)\in
       \operatorname{Reach}_P(M_0)\},\qquad
 \operatorname{Reach}_P(M_0)=\lf(X\mapsto I\cup\Step(X)).
\]
따라서 Exercise 2는 Exercise 1의 도달 상태 중 정상 종료점의 메모리를 정확히 추출한다.
무한 실행은 오른쪽의 중간 상태에는 기여하지만 $\halt$에는 기여하지 않는다.

\paragraph{검산 및 AI 사용 기록.}
이 문서의 설계와 서술에 AI를 사용하였다.
아래는 수식에 대한 검산이며, 유한 실행 예시를 전체 프로그램에 대한 안전성 증명으로
대체하지 않는다. 최종 안전성의 근거는 위의 경우별 simulation과 최소고정점 논증이다.
\begin{itemize}[leftmargin=*,nosep]
 \item \textbf{요구사항 대조:}
 \texttt{project1.pdf}의 두 exercise에 대해 후속 관계, 한걸음 관계,
 실행기 방정식, 요약 연산자 타입, 안전성 조건을 각각 대응시켰다.
 \texttt{project.pdf}의 command 6종과 식 9종이 모두 정의에 나타남을 확인하였다.
 \item \textbf{포인터 대입:}
 $m(p)=\addr(x)$일 때 $W_{p,E}(\{m\})$는 $x$만 갱신한다.
 특히 $E=*p+1$이면 $m[x\mapsto m(x)+1]$을 얻는다.
 이는 주소와 우변을 모두 갱신 전에 계산하는 (5)와 같다.
 \item \textbf{메모리 상관관계:}
 $m_0(x)=0,m_1(x)=10$이면
 $A_{x,x+1}(\{m_0,m_1\})$의 $x$값은 $1,11$이다.
 메모리별 평가를 유지하므로 두 입력 사이의 다른 변수값을 섞지 않는다.
 \item \textbf{입력과 loop:}
 $\Eval_{\mathsf{readInt}+\mathsf{readInt}}(m)=\mathbb Z$이며,
 $\mathsf{readInt}$ 조건은 참/거짓 두 가지가 모두 가능하다.
 $\mathsf{repeat}\ C\ 0$은 본문을 실행하지 않고 후속 위치로 이동한다.
 $\mathsf{repeat}\ (x:=x)\ 1$은 루프 위치에 도달할 수 있지만 정상 종료 결과는 없다.
 \item \textbf{비지역 점프:}
 (4)의 예에서 $x=1,y=0$일 때 $10,20,30,50,70,80,\halt$의 경로를 확인하였다.
 실행기에서도 $H(P_{70})$가 $80$까지 포함하는 residual execution이므로
 $70$의 부분 트리만 실행하고 멈추거나 $50$으로 복귀하지 않는다.
 \item \textbf{고정점 및 타입:}
 reachability에는 상태집합의 최소고정점,
 interpreter에는 command별 함수 환경의 최소고정점을 사용하였다.
 local pair 관계와 label map 변환을 분리하고 $\Graph,\pi,J$의 타입을 명시하였다.
 $\gamma_M(\bot^\sharp)=\emptyset$ 및 $\gamma_M$의 join 보존을 사용하지 않음을
 두 안전성 증명의 각 단계에서 확인하였다.
\end{itemize}

\paragraph{참고 자료.}
\texttt{project1.pdf}; \texttt{project.pdf}.\\
\texttt{lecture slide/slide8.pdf}, PDF pp.~5--9, 17--22, 33--35, 38--44.\\
\texttt{lecture slide/slide7.pdf}, PDF pp.~13--17.\\
\texttt{lecture slide/slide5.pdf}, PDF p.~6.
서식과 저자 표기는 \texttt{hw2.tex}를 따랐다.
Slide 7의 PDF p.~16에서 조건문의 두 번째 가지는 첫 번째 command가 아니라
두 번째 command에 적용되도록 정의하였다.

\end{document}
